% =====================================================================
%  Chapter 10 — Modified Gravity
% =====================================================================

\chapter{Modified Gravity}
\label{ch:modified}

Every argument in this book so far has assumed that the missing mass is
there and that gravity is as we know it. The alternative --- that the
missing mass is not there and that gravity is not as we know it --- has to
be taken seriously, because on the scale of galaxies it works better than
the standard picture.

\section{The problem modified gravity solves}The empirical case for modification is
specific. Under Newtonian gravity, the observed rotation curves of spiral
galaxies require mass-to-light ratios that rise to physically implausible
values in the outskirts, and the central dynamics of dwarf galaxies require
cores rather than cusps. The acceleration scale at which the discrepancy
appears is remarkably consistent across systems: about
$1.2\times 10^{-10}\ \mathrm{m\,s}^{-2}$, independent of the galaxy's
size, mass or surface brightness.

A modification that interpolates from Newtonian to non-Newtonian behaviour at
a fixed acceleration threshold would account for that consistency, and the
family of theories built on that idea --- MOND in its various forms ---
reproduces rotation curves with a single free parameter and no dark halo at
all.

That is a strong result and it should not be dismissed. The point that
follows is that the same success does not extend to the scales at which the
bulk of the evidence lies.

\section{What it cannot reach}Three observations defeat the modified-gravity
programme as it stands, and each is fatal to the alternative rather than
merely inconvenient.

The \emph{Bullet Cluster} is the first and is treated in Chapter~\ref{ch:bullet}.
A theory in which there is no additional mass must still explain the
lensing mass offset from the collisional gas. No current version does.

The \emph{accretion cosmology problem} is the second. In modified gravity,
the growth of structure is altered, and in the leading models it is altered
in the wrong direction: structure grows too fast in the early universe,
producing excessive density fluctuations in the cosmic microwave background
and in the galaxy distribution. The remedy in relativistic models is a
scalar field that restores the growth rate, and that field then behaves as
dark matter in the CMB, which is to say the theory reintroduces the thing it
set out to remove.

The \emph{cluster abundance problem} is the third. Modified gravity without a
halo predicts a certain number of rich clusters per unit volume; the
observed number is far higher, and the deficit is precisely the amount that
additional unseen mass would supply.

\section{The status of the alternative}The fair summary is that modified gravity
is a failed alternative at the scales where the dark matter problem is
actually important, and a successful one at the scales where it is easiest to
study.

That asymmetry is why the field has settled. It is not that the Bullet
Cluster is a perfect experiment --- it is one cluster with its own
uncertainties --- but that three independent arguments fail in the same
direction, and the failures are not correlated in the way that shared
systematics would produce.

There is also a structural argument that is easy to overlook. Modified
gravity is not a prediction; it is a modification with parameters that are
adjusted to fit. Dark matter, whatever it turns out to be, is constrained by
the CMB and by cosmology to a density of 0.12 that no one chose. The
comparison is not between two equally parameterised theories where one is
preferable on aesthetic grounds; it is between a theory with a successful
prediction of an unexplained number and a family of theories that are fitted
to the data in front of them.

That said, the modified-gravity programme has produced physics that the rest
of the field uses. The phenomenological success of the modified-gravity
parameterisations has been a useful forcing function for relativistic
consistency, and the resulting scalar-tensor models are the frameworks in
which the cosmological tensions of Chapter~\ref{ch:tensions} are often
discussed.

\begin{keytakeaways}
\item Modified gravity reproduces galactic rotation curves with one free
parameter and no dark halo, which is a genuine and reproducible success.
\item It fails on the scales where the dark matter problem is actually
important: the Bullet Cluster, the growth of structure, and the cluster
abundance.
\item The asymmetry --- success on galaxies, failure on clusters and
cosmology --- is why the field has settled on dark matter being real.
\item The relativistic versions of these theories introduced the scalar
fields that now carry much of the discussion of cosmological tensions.
\end{keytakeaways}
